\section{Vision as a Perspective：视觉不是一个小领域}

上一章把个人研究轨迹串到 representation，这一章则把“为什么一直做 vision”讲清楚。这里的关键不是为旧学科辩护，而是重新定义 vision 在下一代智能中的位置。

当访谈进入“表征的世界”时，谢赛宁提出一个重要区分：computer vision 不只是 classification、detection、segmentation 这些任务集合，而是一种看待智能的 perspective。它处理的是 continuous, high-dimensional, noisy signals，也就是连续空间、高维、有噪声的信号。这类信号很难被简单 token 化，也不天然带有人类写好的标签。

这解释了为什么 LLM 兴起后，他并不沮丧。LLM 的成功把语言接口、多模态系统和更大规模训练推到前台，反而让 vision 有机会摆脱单个任务，进入“真实世界智能”的大范围问题。真正的危险不是 LLM 太强，而是所有视觉问题都被迫服从语言模型的叙事，最后把视觉退化成 prompt 和 caption 的附庸。

\begin{figure}[H]
\centering
\includegraphics[width=0.94\textwidth]{representation-ladder.png}
\caption{从视觉任务到预测性世界模型：访谈中的能力阶梯。}
\end{figure}

\begin{knowledgebox}{读图：从 task 到 world model 的能力迁移}
图中的 L0 到 L4 不代表官方分级，而是对访谈逻辑的重构。早期视觉任务解决分类、检测、分割；多模态阶段让语言成为接口；再往后，系统必须理解连续事件流、空间结构和 action 后果。终点不是“会看图回答问题”，而是能把视觉流压缩成可用于预测和规划的世界表征。
\end{knowledgebox}

\subsection{层次化表征为什么重要}

谢赛宁反复提到 hierarchical representation。它的直觉是：智能体不可能把世界每一个像素、每一个分子、每一个物理参数都显式记住；它必须学会抽象。抽象不是丢失信息，而是保留和当前任务、未来行动、决策成本有关的信息。

例如房间里的桌子、话筒、光线、声音、纹理都可以被精细建模，但如果目标是继续对话，系统只需要知道话筒能放在桌上、两个人的位置关系、声音是否能被采集等任务相关状态。这个“足够而非全量”的 state，正是 representation learning 和 world model 的交界处。

\begin{importantbox}{表征学习的中心问题}
好的 representation 不是把所有细节重建出来，而是把高维信号压缩成足以支持预测、规划和行动的状态。它要比语言标签更接近物理世界，又要比原始像素更抽象。
\end{importantbox}

\subsection{语言的帮助与污染}

语言作为 interface 极其有用。它让多模态系统可以通过自然语言定义问题、提问和给出答案，也让视觉系统更容易对齐人类目标。但语言也可能成为 shortcut。访谈里谢赛宁用“拐杖”和“鸦片”来形容语言对视觉系统的诱惑：语言让系统看起来更聪明，却可能阻止它训练真正处理连续世界的能力。

\begin{warningbox}{多模态不等于真实理解}
如果一个视觉 benchmark 主要靠语言常识就能答对，那么它并不能证明模型理解了图像、视频或物理世界。语言介入以后，评测必须检查：模型到底用了视觉表征，还是只靠文本先验完成了任务。
\end{warningbox}

\subsection{本章小结}

Vision as a perspective 的核心是：视觉不是被 LLM 接管的小任务，而是构建世界模型所需的底层认知问题。它要求模型处理连续信号、层次抽象、空间结构、事件流和真实行动后果。

